% Why the locked test and VGCure disagree about repair training. Descriptive, written after both were read.
% Numbers come from reports/final/final_primary.json, final_slices.json, final_pairs.json, final_table.json and
% reports/external/vgcure_final_breakdown.json (scripts/vgcure_breakdown.py); checked by tests/test_paper_why.py.
\section{Why the two evaluations disagree}
\label{sec:why}

The two pre-registered evaluations point in opposite directions for repair training: D is
$0.105$ below B on our locked test and $0.056$ above it on the VGCure drawings. They do not ask
the same questions, and splitting the items by task and by answer shows where each number comes
from. This section is descriptive and was written after both results were known; it explains the
numbers and tests nothing. For the same reason, the decision to state the external result in the
abstract and introduction next to the locked-test result was taken after reading both. The wording
rule fixed before the locked test was opened (\texttt{finding-v1}) put only the locked-test statistic
there. The locked-test numbers and their reading are unchanged and the text that rule produced is
kept in the released reports; the hypothesis is still decided by the locked test alone.

\paragraph{On our test, D fails one task.}
On reachability D is level with B ($0.784$ against $0.789$ over the three main locked splits),
although D was never trained on a reachability label. The whole deficit is in the critical-edge
task ($0.529$ against $0.735$). There D answers \emph{true} to $0.792$ of the questions, where
\emph{true} is right for half: it is right on $0.821$ of the items whose answer is true and on
only $0.237$ of those whose answer is false (B: $0.756$ and $0.714$). Over the three main splits
and three seeds it answers both images of a critical-edge twin pair correctly on $228$ of $1{,}800$
pairs, against $856$ for B. D's critical-edge answers follow an answer prior more than the edge a
path depends on. Two development observations fit this (Section~\ref{sec:negative}): providing
the correct edge list as text did not resolve D's critical-edge deficit, although that format was
never seen in training and lowers B's score as well, so it does not rule out a part for perception;
and an auxiliary target added in the one revision we spent moved the prior to \emph{true} on $1\%$
of questions without making the answers right. A possible source, which we did not test: whenever a
repair target removes an edge to break a reachability goal, that edge is by construction critical
for the goal (a single edit can only break reachability by removing such an edge), so in D's
training data every edge that a target removes disconnects something.

\paragraph{On VGCure, D wins on false claims and unreachable pairs.}
VGCure has no critical-edge items, so D's weak task is absent. Its edge claims are half true and
half false. B says that the edge exists to $0.607$ of them and accepts $0.289$ of the false claims;
D says so to $0.513$ and accepts $0.146$. D is right on $0.854$ of the false claims against $0.711$
for B, and slightly behind on the true ones ($0.881$ against $0.925$). On path questions D is ahead
on both classes, most on the unreachable pairs ($0.894$ against $0.795$; reachable pairs $0.860$
against $0.833$). D's edge accuracy is above B's in all ten VGCure sources; its reachability is
above in seven, equal in one and below in two. The same tilt is visible on our own reachability
items, where D is right on $0.797$ of unreachable pairs and $0.771$ of reachable ones (B: $0.775$
and $0.802$).

Why repair training helps here we leave as a hypothesis to test. A repair target names one exact
edge, and a removal must name an edge that is drawn, so repair training may teach the model to
check edges one by one, where B's targets are a single yes or no about paths and B leans towards
\emph{yes} on the foreign drawings. Repair training does not, however, make the model find missing
arrows in the image: the repair prompt lists the edges that may be added, and the generator
requires each of them to be absent from the drawing, so the candidates for an addition come from
the text. Nor can we separate a task effect from a style effect. Compared on the same task,
reachability, D$-$B is $-0.016$ on our style-OOD split and $-0.004$ on the in-distribution split,
against $+0.063$ on VGCure. Our reachability items are balanced between yes and no, so their
accuracy weighs the two classes equally; VGCure's are not ($409$ yes, $91$ no) and are scored by
balanced accuracy, which weighs them equally too. In raw accuracy D is still ahead on VGCure,
$0.866$ against $0.826$ ($+0.040$). Either way, the held-out styles of
our generator do not reproduce the VGCure advantage, and VGCure's drawings, with entity names,
relation labels and other layouts, are a larger shift than those styles.

\paragraph{The parse-then-solve arm.}
E is the strongest arm on our test ($+0.065$ over B) and among the weakest on VGCure ($-0.092$).
Its parser, the \texttt{C\_desc} extraction adapter trained on our drawings, returns no usable graph for $7\%$ of the VGCure images and
drops edges from many others: E accepts only $0.553$ of the true edge claims and $0.511$ of the
reachable pairs, while it rejects false claims well ($0.915$). An exact solver only helps when the
transcription is right, and on foreign drawings it often is not.

\paragraph{What would settle it.}
Runs that we did not make would separate the explanations: critical-edge questions on
externally drawn graphs with known edge lists, which would show whether D's answer prior travels
with it, and edge-existence questions on our locked test, which would show whether D's advantage
on absent edges exists in our own drawing style. A third, a repair variant whose prompt does not
list the edges that may be added, would test whether finding missing arrows in the image matters.
All three fit the released protocol.
